% Draft v3 of §4.2 "Understand the Baseline: Shortest Ping" (2026-11-03).
% Replaces the subsection body of sections/eval-on-error-distance.tex from
% "How a measurement reaches a TG" to the Takeaways; the intro paragraph and
% the Fig. 5-7 block stay as they are. Every number is in
% outputs/analysis/v5/_cross/paper/pro-paper/paper_numbers.s4.md (blocks
% sp-mechanism, sp-taxonomy) or .s5.md (tab:x-cell, x-cell RTT boxes), except
% the 5.4--9.6x RTT ratio, kept from the previous text (see the rednote).
% Percentages, correlations and ratios only: no distance, RTT or count.

% ---- the new schematic, placed before the Fig. 5-7 block --------------------
\begin{figure}[t]
    \centering
    \input{figs/sp-schematic.tex}
    \vspace{-1ex}
    \caption{How S-P's answer relates to a TG and its ingress point $X_{\mathrm{in}}$. The VPs colocated with $X_{\mathrm{in}}$ form the low-RTT tier, and S-P returns one of them. (a) has-$X$: $X_{\mathrm{in}}$ lies in the TG's cell, so S-P's answer does too. (b) no-$X$: $X_{\mathrm{in}}$ lies in another cell, and the VP nearest the TG loses because its path detours through $X_{\mathrm{in}}$.}
    \label{fig:sp-schematic}
\end{figure}

% How a measurement reaches a TG: the ingress point and its low-RTT tier
\noindent \textbf{How a measurement reaches a TG.} Fig.~\ref{fig:network-topology} illustrates the path of a measurement. Traffic from a VP leaves the MNO and enters the TG's network at an interconnect $X$, then travels within that network to the TG. Its propagation distance is thus approximately $D_{prop} = D_{VP\text{-}X} + D_{X\text{-}TG}$, the sum of the great-circle distances from the VP to $X$ and from $X$ to the TG. Which interconnect the traffic enters through is set by the BGP announcements of the TG's network, which change infrequently~\citeneeded{stability of BGP ingress / interconnect selection}. We call the interconnect announced for a TG its \emph{ingress point} and denote it $X_{\mathrm{in}}$. VPs close to one another hand their traffic over at the same interconnect, so for the VPs around $X_{\mathrm{in}}$, $D_{X\text{-}TG}$ is the same and their RTTs differ only by $D_{VP\text{-}X}$. Every interconnect has at least one colocated VP (\S~\ref{sec:dataset-collection}), so the VPs colocated with $X_{\mathrm{in}}$ measure the lowest RTTs to the TG. We call them the \emph{low-RTT tier} (Fig.~\ref{fig:sp-schematic}). \emph{S-P returns a VP of this tier, so its prediction lands near $X_{\mathrm{in}}$ rather than near the TG, and its error distance is approximately $D_{X\text{-}TG}$.}

% S-P answers at the ingress point: three observations
\noindent \textbf{S-P answers at the ingress point.} We do not record the ingress point of a TG during measurement, so we use the TG's nearest interconnect $X$ in its place\footnote{We pick $X$ from the TG's location alone, before looking at S-P's answer, so each test below could fail.}. We show three observations that follow from the mechanism above. First, S-P's answer lands in the cell of $X$ for 85.1\% of the TGs, against 6.7\% for a random VP among those measuring the TG, and S-P's VP shares its nearest interconnect with the TG for 81.9\% of the TGs. Second, the VP nearest the TG rarely measures the lowest RTT. This holds for only 13\% of the has-$X$ TGs (defined below) and for none of the no-$X$ TGs, where the nearest VP measures a median of 1.8$\times$ the lowest RTT. Third, S-P's error follows the TG's distance to its nearest interconnect much more closely than its distance to the nearest VP, as shown in Fig.~\ref{fig:sp-scatter}. Spearman's $\rho$ is 0.86 against the former (Fig.~\ref{fig:sp-x-scatter}) and 0.38 against the latter (Fig.~\ref{fig:sp-vp-scatter}), and 69\% of the TGs lie within 2$\times$ of the diagonal in Fig.~\ref{fig:sp-x-scatter}, against 35\% in Fig.~\ref{fig:sp-vp-scatter}. The nearest interconnect is not always the ingress point. For 4.5\% of the TGs, the lowest RTT is shorter than the propagation time of any path through $X$, so their traffic must enter the network elsewhere. \emph{S-P's prediction is set by where the TG's traffic enters the network, and the VP nearest the TG has little say in it.}

% S-P is correct only when the TG's cell holds its ingress point
\noindent \textbf{S-P is correct only when the TG's cell holds its ingress point.} The answer space (\S~\ref{sec:answer-space}) lets us check this directly. We call a TG \emph{has-$X$} if its own cell holds an interconnect, and \emph{no-$X$} otherwise. Pooled over the three networks, 53\% of the TGs are has-$X$ and 47\% are no-$X$. S-P's answer lands in a cell that holds an interconnect for every TG, against 67\% for a random VP. A no-$X$ TG's cell holds no interconnect, so S-P cannot classify it correctly, whichever interconnect its traffic enters through (Fig.~\ref{fig:sp-schematic}b). A has-$X$ TG is classified correctly when its traffic enters at the interconnect in its own cell (Fig.~\ref{fig:sp-schematic}a). Fig.~\ref{fig:sp-accuracy-by-x} confirms both cases. S-P classifies 93\% of the has-$X$ TGs correctly and none of the no-$X$ TGs, and 88\% of its wrong no-$X$ answers land in a cell adjacent to the TG's. The has-$X$ label alone predicts S-P's outcome for 96\% of the TGs. Per network, S-P's accuracy is close to the network's has-$X$ share (Table~\ref{tab:x-cell}: 64\% vs.\ 75\% in AS-A, 42\% vs.\ 42\% in AS-B, 43\% vs.\ 43\% in AS-C). VP proximity does not separate the two groups, since the nearest VP of a no-$X$ TG is about as close as that of a has-$X$ TG (a median of 1.1$\times$ farther). Their minimum RTTs do separate them, as shown in Fig.~\ref{fig:boxplot-rtt-of-x-cell-by-asn}. Per network, the median minimum RTT of the no-$X$ TGs is 5.4--9.6$\times$ that of the has-$X$ TGs\rednote{5.4--9.6$\times$ is not in paper\_numbers yet; recompute from x\_cell\_rtt.csv}, and a single RTT level splits the groups: every no-$X$ TG lies above it and 92\% of the has-$X$ TGs below it (the dotted line). \emph{S-P is correct only when the TG's cell holds its ingress point, so the has-$X$ label explains its accuracy.}

% RTT discriminative power: the has-X misses
\noindent \textbf{RTT discriminative power is the second driver of accuracy.} S-P relies on the minimum RTT to single out the low-RTT tier, and within the tier the VP nearest the ingress point. When the RTTs of the candidate VPs are too similar to tell apart, S-P can return a VP that is not the closest to the ingress point. This accounts for the 7\% of has-$X$ TGs that S-P classifies wrongly (Fig.~\ref{fig:sp-accuracy-by-x}). All of them are in AS-A, and they come in two kinds. (1) \emph{Near-ties} (3.0\% of the has-$X$ TGs). VPs within 100~km of the same $X$ differ by only about 1~ms in propagation round-trip time at the speed of the Internet, and a slightly different access delay at the TG can outweigh that difference. When the tied VPs sit in different cells, S-P picks among them almost at random. The same effect spreads S-P's error at the lower left of Fig.~\ref{fig:sp-x-scatter}, where it is not strictly proportional to the distance to $X$. (2) \emph{RTT that ignores distance} (3.7\%). Every VP measures a heavily inflated minimum RTT to these TGs, at least 6.3$\times$ the largest of any correctly classified has-$X$ TG, due to a deliberate routing policy we confirmed internally. Their RTT no longer tracks distance. The per-TG Spearman's $\rho$ between RTT and VP distance is at most 0.20 for them, against at least 0.81 for every correctly classified has-$X$ TG. S-P's predictions then fall far from these TGs, and they form the outliers at the top of Fig.~\ref{fig:sp-x-scatter}. \emph{When RTTs lose their discriminative power, S-P misses even the TGs whose cell holds the ingress point.}

% In summary, here is how S-P works in our setting and why its accuracy is what it is.
\noindent \textbf{Takeaways.} \textit{In an operator's setting, every interconnect has a colocated VP, so the low-RTT tier of VPs sits at the TG's ingress point. Shortest Ping's prediction therefore lands in the cell of the ingress point, and it is correct when the ingress point lies in the TG's serving region. Similar or inflated minimum RTTs among VPs can still confuse Shortest Ping and lead to a wrong prediction.}
